% Conversão de chapter-06-text-h-0110; origem RELATORIO_ATUAL.odt.
\hypertarget{__RefHeading___Toc4837_3118176633}{}
\refstepcounter{subsection}
\subsection*{5.2.3 Arquitetura utilizada}
\label{ages-iv:text-h-0110}
\addcontentsline{toc}{subsection}{5.2.3 Arquitetura utilizada}

% SOURCE text-p-1347
A arquitetura da solução foi projetada com foco em desacoplamento, modularidade e manutenibilidade entre as camadas mobile e backend. O frontend consiste em um aplicativo mobile desenvolvido em React Native\space\cite{ages-iv-react-native} com Expo\space\cite{ages-iv-expo}, estruturado em uma arquitetura orientada a módulos funcionais com gerenciamento de estado de servidor via TanStack Query\space\cite{ages-iv-tanstack-query} e preparado para distribuição em ambiente iOS e Android. O backend, construído em Python\space\cite{ages-iv-python} com FastAPI\space\cite{ages-iv-fastapi} e conteinerizado via Docker\space\cite{ages-iv-docker}, expõe uma API REST assíncrona organizada em camadas, centralizando as regras operacionais da plataforma, validação estrita de contratos de dados com Pydantic\space\cite{ages-iv-pydantic} e integração com APIs externas de busca musical.\par

A equipe de AGES III implantou a API em uma instância Amazon EC2, com o Caddy recebendo as conexões HTTPS e encaminhando as requisições ao FastAPI. O Amazon Route 53 resolve o domínio da API. A API acessa o PostgreSQL 16 hospedado em uma instância privada do Amazon RDS\space\cite{ages-iii-rds}. As fotos e os vídeos são armazenados em um bucket privado do Amazon S3\space\cite{ages-iii-s3}. A API gera URLs pré-assinadas para que o aplicativo envie os arquivos diretamente ao bucket, e o Amazon CloudFront\space\cite{ages-iii-cloudfront} distribui as mídias a partir desse bucket.\par

Complementando a infraestrutura, o ecossistema utiliza esteiras automatizadas de CI/CD via GitHub Actions\space\cite{ages-iv-github-actions}, integradas a análises estáticas de qualidade de código (SonarCloud\space\cite{ages-iv-sonar}) e execução de testes automatizados com Pytest\space\cite{ages-iv-pytest} e Jest\space\cite{ages-iv-jest}. A esteira de implantação gera a imagem Docker para ARM64 e a publica no GitHub Container Registry. A EC2 baixa a imagem e executa o backend com Docker Compose.\par

% SOURCE text-a-0150 (âncora ODF vazia; destino preservado sem texto visível)
% Origem text-a-0150: https://tools.ages.pucrs.br/2026-2/2jk-4jk/recorda/wiki/-/wikis/arquitetura

% FIGURE draw-image-0026 / CAPTION text-p-1350
\memorialfigure{imagens/user-supplied/recorda-arquitetura.png}{\textbf{Figura 25 –} Diagrama de Arquitetura - AGES IV}{ages-iv:draw-image-0026}{Fonte: \url{https://tools.ages.pucrs.br/2026-2/2jk-4jk/recorda/wiki/-/wikis/arquitetura}}{0.98}

% SOURCE text-p-1362 / text-a-0151
Mais informações sobre a arquitetura estão disponíveis em: \url{https://tools.ages.pucrs.br/2026-2/2jk-4jk/recorda/wiki/-/wikis/arquitetura}\par
